\chapter{The Hallucination Problem, and Why \texorpdfstring{\textsc{CAEM}}{CAEM} Exists}
\label{ch:problem}

Before we touch a single component, we need to feel the problem clearly. Every
design choice in \caem{}, namely the memory, the nine signals, the calibrated
gate, the self-improvement loop, and the rollback, is a direct response to a
specific failure of ordinary language models. Once we understand that failure,
the architecture stops looking like a pile of clever tricks and starts looking
like the \emph{only} sensible response. That is the goal of this chapter.

\begin{intuition}
A language model is a fluent guesser. It is extraordinarily good at producing
text that \emph{sounds} right. The trouble is that ``sounds right'' and ``is
right'' are different properties, the model has no built-in way to tell them
apart, and, worse, no way to signal which is which to a reader. A confident lie
and a confident truth come out of the model wearing the same clothes. \caem{}
exists to put a disciplined system around the guesser, so that the answers it
commits to are trustworthy, it knows how much to trust each one, and it gets
better over time without forgetting what it already knew.
\end{intuition}

\section{What a hallucination actually is}

A \emph{hallucination} is a model output that is fluent and confident but
factually wrong. The word covers a range of causes, from a retrieval system
that missed, to a genuine gap in the model's knowledge, to outright
confabulation under uncertainty. The symptom is always the same: the model
states something untrue as though it were true.

The danger is not the wrongness by itself. We tolerate wrong answers from
search engines, from textbooks, and from each other. The danger is the
\textbf{calibration mismatch}: the model expresses the same high confidence
whether it is right or wrong, so a reader with no independent way to check
cannot tell which answers to trust. A false answer phrased like a true one is
harder to catch than a false answer that sounds unsure. In a casual chat that
is an annoyance. In a clinical, legal, or financial setting it is a hazard.

\section{Scale does not fix it}

The first instinct of the field was that hallucination is a small-model problem
that scale will erase. The evidence says otherwise.

\begin{itemize}[leftmargin=1.4em]
  \item On \textbf{TruthfulQA}, a benchmark deliberately built around questions
  where humans hold confident misconceptions, the largest frontier model
  evaluated reached only \textbf{58\%} on the truthfulness axis and \textbf{25\%}
  on the combined truthful-and-informative axis, against a human baseline of
  \textbf{94\%} on both~\cite{lin-etal-2022-truthfulqa}.
  \item The same benchmark documented \textbf{inverse scaling}: larger models
  were \emph{more}, not less, prone to echo human falsehoods. Bigger made the
  result worse on exactly the questions that matter~\cite{lin-etal-2022-truthfulqa}.
  \item On \textbf{long-form generation}, the FactScore atomic-fact decomposition
  protocol found that only \textbf{58\%} of decomposed atomic claims in
  chat-style biographical generation were supported by evidence; the rest ranged
  from unsupported speculation to outright fabrication~\cite{min-etal-2023-factscore}.
  \item In \textbf{clinical question answering}, frontier models hallucinated in
  \textbf{69--77\%} of responses, an order of magnitude above any threshold a
  responsible deployment would tolerate~\cite{alkaissi2023artificial}.
\end{itemize}

\begin{pitfall}[A number that appeared on the poster]
An earlier draft of the poster quoted a ``3 to 91 percent'' deployment-domain
range. That specific range is \emph{not} in the thesis, and we do not use it
here. The grounded figures are the four above. If the 3--91\% number comes up,
it should be treated as a loose summary rather than a thesis claim. This book
uses only what the report actually establishes.
\end{pitfall}

The survey literature draws the conclusion bluntly: hallucination is
``a structural property of stateless generation under uncertainty rather than a
defect to be trained away''~\cite{ji2023survey,huang-etal-2023-large}. That
sentence rewards a second reading. It says the problem is not in the
\emph{weights}; it is in the \emph{way the model is used}. A bigger guesser is
still a guesser. \textbf{Any system that aims to reduce hallucination must
change the structure, not just train harder.} That single sentence is the seed
of the entire thesis.

\section{The three structural flaws}
\label{sec:three-flaws}

The thesis pins the structural problem to three specific shortcomings of the
standard inference paradigm. These three words, \textbf{stateless},
\textbf{overconfident}, and \textbf{static}, are worth committing to memory,
because every part of \caem{} maps back to one of them.

\subsection*{Flaw 1: Stateless}
Standard inference processes every query independently. No representation of
past successful answers is kept across sessions, and the same question asked
twice can produce two different answers. A reasoning chain that produced a
correct answer this morning is discarded the moment the answer is delivered.
This wastes computation on questions the system has effectively already solved,
and it prevents any accumulation of domain knowledge through use. The model
never gets to say ``I have seen this before, and last time I got it right.''

\subsection*{Flaw 2: Overconfident (miscalibrated)}
Modern neural networks are systematically overconfident. Calibration studies
report that both classifiers and language models assign high probability to
incorrect outputs at rates that exceed their actual accuracy by \textbf{ten
percentage points or more}~\cite{pmlr-v70-guo17a, kadavath2022language}. The
obvious fix, reading off the model's own confidence, does not work either:
single-signal correctness predictors achieve area-under-the-ROC scores between
only \textbf{0.50 and 0.60}, barely above random chance~\cite{xiong2023can}.
The consequence is the one we opened with: a user cannot reliably tell which
answers are likely wrong from confidence alone.

\subsection*{Flaw 3: Static}
Once trained, a deployed model neither learns from its mistakes nor adapts to
the institution using it, because naive online fine-tuning incurs
\textbf{catastrophic forgetting}: it improves on new material by erasing the
broad capabilities the base training instilled~\cite{doi:10.1073/pnas.1611835114}.
The implication is stark. Verification effort spent today produces no
compounding benefit tomorrow. A thousand answers verified by humans this morning
will not change a single answer the model gives this afternoon.

\subsection*{Why the three are a knot, not a list}

The crucial observation is that these flaws \textbf{reinforce one another}. They
are not three independent bugs that could be fixed one at a time; they form a
cycle.

\begin{itemize}[leftmargin=1.4em]
  \item Without \textbf{memory}, the model cannot reuse verified knowledge.
  \item Without \textbf{calibrated confidence}, the system cannot recognise when
  memory is missing and worth consulting, so even a perfect memory would be
  under-used.
  \item Without \textbf{self-improvement}, the verifications that humans or
  external systems do produce decay into one-shot interventions.
\end{itemize}

\begin{figure}[t]
\centering
\begin{tikzpicture}[
  font=\small\sffamily,
  flaw/.style={rectangle, rounded corners=6pt, draw=#1!75!black, fill=#1!10,
    line width=1.1pt, align=center, text width=3.5cm, minimum height=1.55cm,
    inner sep=5pt, drop shadow={shadow xshift=0.4ex,shadow yshift=-0.4ex,
    fill=black!12}},
  rel/.style={-{Stealth[length=3mm]}, line width=1.1pt, draw=black!50},
]
  \node[flaw=cbIntuition] (state) at (90:3.1) {\textbf{Stateless}\\[2pt]
    {\footnotesize no verified answer is retained}};
  \node[flaw=cbPitfall]   (conf)  at (210:3.1) {\textbf{Overconfident}\\[2pt]
    {\footnotesize cannot tell right from wrong}};
  \node[flaw=cbExample]   (static)at (330:3.1) {\textbf{Static}\\[2pt]
    {\footnotesize verification never compounds}};
  \node[circle, draw=black!25, fill=black!3, line width=0.8pt,
        text width=1.7cm, align=center, font=\scriptsize] (c) at (0,0)
    {the\\ reinforcing\\ cycle};

  \draw[rel] (state) to[bend left=20]
    node[midway,right,text width=2.5cm,align=left,font=\scriptsize]
    {no memory, so nothing to calibrate about} (conf);
  \draw[rel] (conf) to[bend left=20]
    node[midway,below,text width=2.8cm,align=center,font=\scriptsize]
    {cannot spot gaps, so cannot target learning} (static);
  \draw[rel] (static) to[bend left=20]
    node[midway,left,text width=2.5cm,align=right,font=\scriptsize]
    {no learning, so verified knowledge never persists} (state);
\end{tikzpicture}
\caption[The three-flaw cycle]{The three structural flaws form a reinforcing
cycle, not an independent list. Fixing any one in isolation leaves the other two
to drag it back. The thesis of \caem{} is that they must be attacked
\emph{simultaneously}. (Schematic.)}
\label{fig:three-flaw-cycle}
\end{figure}

This is why fixing one flaw alone fails: each unsolved flaw pulls the solved one
back toward the broken state. \Cref{fig:three-flaw-cycle} draws the cycle. The
whole architecture of \caem{} is the claim that the cycle can be broken along
all three axes at once: memory for statelessness, a calibrated verifier for
overconfidence, and a guarded self-improvement loop for staticness.

\section{Why existing fixes do not close the gap}
\label{sec:prior-fail}

Three families of prior systems each attack hallucination, and each has a
structural limitation that prevents it from closing the gap alone.

\begin{description}[leftmargin=0pt, labelsep=0.5em, style=nextline]
  \item[Retrieval-augmented generation (RAG).] RAG grounds the model in
  retrieved evidence~\cite{NEURIPS2020_6b493230}, which raises factual recall
  \emph{when retrieval succeeds}. It is silent, however, when retrieval misses
  or returns evidence the model summarises incorrectly. The output is still a
  single-pass generation whose quality is capped by retrieval quality, and the
  system keeps \emph{no record of which past answers were correct}. RAG does not
  touch statelessness.

  \item[Prompting interventions.] Chain-of-thought reasoning and
  self-consistency over sampled chains~\cite{wei2022chain, wang2023selfconsistency}
  expose internal reasoning and exploit agreement across resampled answers. They
  provide \emph{no persistence}, however: each session begins with no memory of
  prior verified answers, and the same question tomorrow triggers the same
  reasoning effort as today. They make a single answer better; they do not make
  the \emph{system} better.

  \item[Post-generation verification.] Methods such as SelfCheckGPT and semantic
  entropy~\cite{manakul-etal-2023-selfcheckgpt, farquhar2024semantic} detect
  hallucinations after they occur and support refusal or revision. The verdict
  \emph{does not feed back}, however, into the model's weights or any memory, so
  the identical hallucination recurs in the next session. Verification without a
  feedback loop is a one-shot intervention.
\end{description}

\begin{intuition}[The pattern across all three]
RAG fixes grounding but not statelessness. Prompting fixes single-answer quality
but not persistence. Verification detects errors but never feeds the fix back.
Each family solves a different \emph{slice} of the problem and leaves the rest
untouched, which is exactly why stacking ``a stronger retriever, longer chains,
a sharper verifier'' yields diminishing returns. None of them closes the loop.
\end{intuition}

The thesis states the implication directly: hallucination reduction at
deployment-relevant rates ``is unlikely to come from incremental improvement of
any one family.'' What is missing is \textbf{a closed loop}, a system in which
verification feeds memory, memory feeds routing, routing feeds generation, and
improvement compounds across cycles instead of evaporating at session
boundaries. Arguing for that system, and building it, is the work of the thesis.

\section{The architectural opportunity}
\label{sec:arch-opportunity}

If the three flaws are a knot, the opportunity is that three capabilities, taken
\emph{together}, untie it. \caem{} commits to exactly three, each aimed at one
flaw.

\begin{enumerate}[leftmargin=1.6em]
  \item \textbf{A verified episodic memory with a four-outcome admission rule},
  finer than a single binary threshold. Outputs are sorted into store, defer,
  abstain, or discard, so borderline answers wait in a deferred buffer
  \emph{without contaminating the high-precision pool}. This attacks
  statelessness.
  \item \textbf{Confidence-aware routing with a retrieval-utility decision}. A
  calibrated pre-generation confidence and a learned classifier decide
  \emph{how hard to work} and \emph{whether retrieval will even help} on each
  query. This attacks overconfidence.
  \item \textbf{A constrained self-improvement loop with retention rollback}.
  The model re-trains on its own verified answers each cycle, but a guard
  reverts any update that erodes prior capability. This attacks staticness.
\end{enumerate}

Taken together, these convert the two statelessness problems into two
\emph{stateful} mechanisms: the verified store absorbs recurring queries at
lookup cost, and the self-improvement loop folds verified retrieval-augmented
reasoning into the model's own parametric capacity. The three serving tiers,
memory, zero-shot, and retrieval, stop being fixed cost categories and become a
\textbf{self-organising compute ladder} under the cycle loop.

\begin{precise}[The data-purity precondition, or why an imperfect verifier is enough]
The natural objection is this: if the verifier is imperfect, will it not admit
wrong answers and poison the memory? The thesis answers with a base-rate
argument. Under reasonable conditions on the verifier's balanced accuracy and
the base generation accuracy, \textbf{the purity of the stored pool can exceed
the verifier's own accuracy}, because base-rate dynamics favour the gate
whenever the verifier is more likely to admit a correct answer than an incorrect
one. The precondition is mild: the verifier's balanced accuracy need only exceed
one half on the distribution of claims it actually scores. We prove and test
this in \Cref{ch:theory}. For now, the intuition to hold is that a perfect
verifier is not required, only one that is better than a coin flip on the claims
it sees. On the realised system this precondition is met with room to spare: the
deployed verifier's balanced accuracy on the calibration-matched distribution is
about 0.67, and the stored pool's realised purity sits near 0.74, above the
verifier's own accuracy exactly as the identity predicts (the proof is in
\Cref{ch:theory} and the empirical receipt in the evaluation chapter).
\end{precise}

\section{The scope of the claim}

So the thesis does not get to claim that \caem{} solves hallucination. It claims
something sharper and falsifiable. The intervention is bounded so that it can be
argued empirically within an undergraduate research budget.

\begin{itemize}[leftmargin=1.4em]
  \item \textbf{One small frozen generator}: Qwen-2.5-3B-Instruct, a
  three-billion-parameter open-weight instruction-tuned decoder. The backbone is
  \emph{frozen}; only a low-rank adapter on the attention and feed-forward
  projections trains, at roughly \textbf{2\%} of the parameter count.
  \item \textbf{One consumer GPU}: the whole system fits on a single
  thirty-two-gigabyte graphics processor under sixteen-bit precision, so the
  result is reproducible without specialised infrastructure.
  \item \textbf{Five benchmarks}: a three-benchmark \emph{training panel}
  (FEVER, TriviaQA, CommonsenseQA) and a two-benchmark \emph{transfer panel}
  (TruthfulQA, StrategyQA), with multitask language understanding (MMLU) held
  out as an out-of-distribution retention control rather than an optimisation
  target.
  \item \textbf{A fixed open generator}, with open verification machinery and no
  external commercial calls, so the architectural claim is separable from any
  pre-training advantage and reproducible end-to-end.
\end{itemize}

\begin{defence}[The one-sentence framing of the whole thesis]
``Hallucination is a structural property of stateless, overconfident, static
generation, not a defect that scale removes. \caem{} attacks all three at once
on a fixed small model: a verified episodic memory for statelessness, a
calibrated nine-signal verifier for overconfidence, and a guarded
self-improvement loop for staticness; and we show the gains compound across
cycles while a retention guard prevents regression.''
\end{defence}

\section{What the rest of the book does}

We now have the \emph{why}. The rest of the book is the \emph{how}, built
bottom-up.

\begin{itemize}[leftmargin=1.4em]
  \item \Cref{ch:bigpicture} gives the whole machine in one picture, the
  eight-stage per-query pipeline and the cycle loop that wraps it, so we have a
  map before zooming in.
  \item \Cref{ch:prereq} fills in the concepts the later chapters lean on
  (embeddings, calibration, natural-language inference, low-rank adaptation, and
  similarity search), so the book is self-contained.
  \item Each stage then gets its own chapter, grounded in the live code: what it
  does, the exact mechanism, a worked example, the traps, and how to defend it.
\end{itemize}

It helps to keep \Cref{fig:three-flaw-cycle} in mind throughout. Whenever a
component seems elaborate, we can ask which of the three flaws it is breaking the
cycle on. The answer is always one of the three.
